Fix \(P\in\mathcal P_{d,\epsilon,M,\sigma}\) and normalize by
\[
 V:=Y/M,quad \nu_{ak}:=\mu_{ak}/M,quad
 s_{ak}:=P(X=k,A=a),quad z_{ak}:=E[V\mathbf 1\{X=k,A=a\}],
\]
and \(\varphi_k:=p_k(\nu_{1k}-\nu_{0k})\).  Then
\[
 \tau(P)/M=\sum_{k=1}^d\varphi_k,qquad
 z_{ak}=s_{ak}\nu_{ak},qquad |\nu_{ak}|\le1/2,qquad s_{ak}\ge\epsilon p_k.
\tag{1}
\]
The last two relations follow from mean normalization and overlap.  The new moment assumption gives directly
\[
 \operatorname{Var}(V\mid A=a,X=k)le1,
\tag{2}
\]
which is the only tail control used by either the light or heavy part.

Let \(T_K\) be the Chebyshev polynomial of the first kind and, with continuous extension at zero, define
\[
 H_K(x):=\frac{1-T_K(1-2x)}{2K^2},qquad
 E_K(x):=\frac{H_K(x)}x,qquad
 G_K(x):=\frac{1-E_K(x)}x.
\]
Expanding the shifted Chebyshev polynomial yields
\[
 G_K(x)=\sum_{j=0}^{K-2}(-1)^j
 \frac{2^{2j+3}}{K(K+j+2)}\binom{K+j+2}{2j+4}x^j
 =\sum_{j=0}^{K-2}g_{Kj}x^j.
\tag{3}
\]
For \(0\le x\le1\), \(|T_K(1-2x)|\le1\), so
\[
 0\le x\{1-xG_K(x)\}=H_K(x)\le K^{-2}.
\tag{4}
\]
If \(p_k\le B\), equations (1) and (4) imply
\[
 \left|p_k\nu_{ak}-\frac{p_kz_{ak}}B G_K(s_{ak}/B)\right|
 \le\frac{B}{2\epsilon K^2}.
\tag{5}
\]
Lemma~\(\mathrm{lem:linear\text{-}mark\text{-}factorial\text{-}covariance}\) gives \(E[U_{k,a,j}]=z_{ak}p_ks_{ak}^j\); hence \(\widehat P_k\) is unbiased for
\[
 \frac{p_kz_{1k}}B G_K(s_{1k}/B)-
 \frac{p_kz_{0k}}B G_K(s_{0k}/B).
\tag{6}
\]

For the independent pilot block put
\[
 B_-:=\frac{\lambda_0L}{2m_0},\qquad B_+:=\frac{2\lambda_0L}{m_0}
\]
and let
\[
 \mathcal A:=\{\widehat{\mathcal H}\subseteq\{k:p_k\ge B_-\},\quad
 \widehat{\mathcal L}\subseteq\{k:p_k\le B_+\}\}.
\]
The balanced split and \(b_0=16\lambda_0\) give \(B_+\le B/4\) for all sufficiently large \(n\).  The binomial Chernoff inequalities and a union bound yield
\[
 Q_P^{(n)}(\mathcal A^c)\le2d\exp(-c\lambda_0L)\le Cn^{-6}
\tag{7}
\]
uniformly for \(d\le\rho_\epsilon nL\), after taking \(\rho_\epsilon\le1\).

Condition on the pilot and suppose \(\mathcal A\) holds.  Then the light set is deterministic relative to \(I_1\), is contained in \(\{k:p_k\le B/4\}\), and, after increasing \(N_\epsilon\),
\[
 K\ge2,qquad4(K+2)^2\le m_1,qquad4(K+2)/m_1\le3B/4.
\tag{8}
\]
Equations (5)--(6) and the covariance lemma therefore give
\[
 \begin{split}
 E\left[\left\{\sum_{k\in\widehat{\mathcal L}}\widehat P_k-
 \sum_{k\in\widehat{\mathcal L}}\varphi_k\right\}^2\middle|I_0\right]
 \le C_\epsilon\left[\frac1{m_1}+6^{2K}\left\{dB^2+
 \frac{d^2K^2B^2}{m_1}\right\}+\frac{d^2B^2}{K^4}\right].
\end{split}
\tag{9}
\]
Let \(\kappa:=2\alpha_0\log6\le1/32\).  Since \(6^{2K}\le n^\kappa\), \(B\le CL/n\), and \(K\le CL\),
\[
 n^\kappa\frac{dL^2}{n^2}\le C\left\{\frac1n+\frac{d^2}{n^2L^2}\right\},
 \qquad n^\kappa\frac{d^2L^4}{n^3}\le C\frac{d^2}{n^2L^2},
\tag{10}
\]
while \(dB/K^2\le Cd/(nL)\).  Thus
\[
 E\left[\left\{\sum_{k\in\widehat{\mathcal L}}\widehat P_k-
 \sum_{k\in\widehat{\mathcal L}}\varphi_k\right\}^2\middle|I_0\right]
 \le C_\epsilon\left\{\frac1n+\frac{d^2}{n^2L^2}\right\}.
\tag{11}
\]

For the heavy cells, conditional on the pilot and on the category and treatment indicators in \(I_1\), (2) gives
\[
 \operatorname{Var}\left(\sum_{k\in\widehat{\mathcal H}}\widehat h_k
 \middle|(X_i,A_i)_{i\in I_1}\right)
 \le\frac1{m_1^2}\sum_{k\in\widehat{\mathcal H}}N_k^2
 \left\{\frac{\mathbf1\{N_{1k}>0\}}{N_{1k}}+
 \frac{\mathbf1\{N_{0k}>0\}}{N_{0k}}\right\}.
\tag{12}
\]
For \(D\sim\operatorname{Binomial}(t,q)\), the identity
\[
 E[(D+1)^{-1}]=\frac{1-(1-q)^{t+1}}{(t+1)q}
\]
and \(D^{-1}\mathbf1\{D>0\}\le2(D+1)^{-1}\) show that the expectation of (12) is at most \(C_\epsilon/m_1\).

The conditional outcome mean of the heavy sum equals
\[
 \frac1{m_1}\sum_{k\in\widehat{\mathcal H}}N_k(\nu_{1k}-\nu_{0k})-R_1+R_0,
 \qquad R_a:=\frac1{m_1}\sum_{k\in\widehat{\mathcal H}}N_k\nu_{ak}\mathbf1\{N_{ak}=0\}.
\tag{13}
\]
The first term is an empirical average of a function bounded by one, is unbiased for \(\sum_{k\in\widehat{\mathcal H}}\varphi_k\), and has variance at most \(1/m_1\).  Exact multinomial factorial moments, \(|\nu_{ak}|\le1/2\), and \(s_{ak}\ge\epsilon p_k\) give
\[
 E[R_a^2]\le C_\epsilon\left[\frac1{m_1}+
 \left\{\sum_{k\in\widehat{\mathcal H}}p_ke^{-c_\epsilon m_1p_k}\right\}^2\right].
\tag{14}
\]
On \(\mathcal A\), each selected heavy cell has \(p_k\ge B_-\ge cL/n\), whence
\[
 E[R_a^2]\le C_\epsilon\left\{\frac1n+\frac{d^2}{n^2L^2}\right\}.
\tag{15}
\]
Combining (12)--(15) proves the heavy analogue of (11).

The selected heavy and light sets partition all cells.  If \(Z\) is the normalized unclipped sum, the two bounds, (1), and \((x+y)^2\le2x^2+2y^2\) yield
\[
 E[\mathbf1_{\mathcal A}\{Z-\tau(P)/M\}^2]
 \le C_\epsilon\left\{\frac1n+\frac{d^2}{n^2L^2}\right\}.
\tag{16}
\]
Projection onto \([-1,1]\) cannot increase loss because \(\tau(P)/M\in[-1,1]\), and its squared loss is at most four on \(\mathcal A^c\).  Equations (7) and (16) prove the claimed rate when \(n\ge N_\epsilon\) and \(d\le\rho_\epsilon nL\).  If \(d>\rho_\epsilon nL\), the zero fallback has loss at most \(M^2\) and \(\min\{1,d^2/(n^2L^2)\}\ge\rho_\epsilon^2\); if \(n<N_\epsilon\), the same conclusion follows from \(1\le N_\epsilon/n\).  Thus the statement holds throughout \(1\le d\le c_\epsilon n^2/L\).

All ratios are multiplied by their positive-denominator indicators.  After cell aggregation the unnormalized ordered sum in \(U_{k,a,j}\) is
\[
 \left(\sum_{i\in I_1}V_i\mathbf1\{X_i=k,A_i=a\}\right)
 (N_{ak}^{(1)}-1)_j(N_k^{(1)}-j-1),
\]
with value zero when too few distinct indices exist.  Hence the estimator is measurable and total, takes \(O(dK^2)\) arithmetic operations, and uses no moment of \(Y\) beyond the conditional second moment in (2).